\section{NN-S017: the counting cubic is an integral wedge of motivic spheres}

\paragraph{Statement and precise scope.}
Retain the actual integral verifier family of NN-C027 and NN-S016,
with $v,h\ge1$:
\[
 X_\Phi=\{t(t-1)=\textstyle\sum_{i=1}^h y_i f_i(x)\},\qquad
 Y_\Phi=V(f_1,\ldots,f_h)=\coprod_{a\in\mathcal S}\operatorname{Spec}\mathbb Z,
 \qquad N=|\mathcal S|.
\]
All $v$ Boolean equations $x_j^2-x_j$ belong to the list. Point
$X_\Phi$ by $x=0,t=0,y=0$; this point exists even when the zero
assignment is not a satisfying assignment. In the pointed unstable
Morel--Voevodsky homotopy category over $\mathbb Z$ there is an isomorphism
\begin{equation}\label{eq:S017-motivic-wedge}
 (X_\Phi,0)\simeq_{\mathbb A^1}
       \bigvee_{a\in\mathcal S}(\mathbb P^1,\infty)^{\wedge v}.
\end{equation}
The empty wedge means the pointed base scheme. The same statement holds
after every base change, with the same assignment index set and normal
framing. In particular this identifies the unstable motivic homotopy
type, not only its realizations, its Grothendieck class, or its ordinary
cohomology.

For every perfect field $k$, its integral geometric motive consequently is
\begin{equation}\label{eq:S017-motive}
 \mathrm M(X_{\Phi,k})\simeq
       \mathbb Z\oplus\mathbb Z(v)[2v]^{\oplus N}.
\end{equation}
For $q,n\ge0$, the entire higher Chow groups are
\begin{equation}\label{eq:S017-higher-chow}
 \operatorname{CH}^{q}(X_{\Phi,k},n)\simeq
 \operatorname{CH}^{q}(k,n)\oplus
       \operatorname{CH}^{q-v}(k,n)^{\oplus N},
\end{equation}
where negative codimension means zero. More generally, for every represented
bigraded motivic cohomology theory $E$ over the chosen base,
\[
 E^{p,q}(X_\Phi)\simeq E^{p,q}(\operatorname{Spec}\mathbb Z)
       \oplus E^{p-2v,q-v}(\operatorname{Spec}\mathbb Z)^{\oplus N},
\]
with the corresponding base changed formula. These are isomorphisms of
groups or modules; no multiplicative assertion is hidden in this formula.

On the specified verifier encoding, the number of nontrivial motivic
sphere summands, equivalently the displayed Tate multiplicity, is
\#P-complete. Recognizing $\mathbb A^1$-contractibility, triviality of the
reduced motive, or the absence of a nontrivial sphere is coNP-complete.
The motive statement is asserted over a fixed perfect field; the
contractibility statement is valid over $\mathbb Z$ and over any fixed
field. These are complexity claims for this effectively constructed
family, not decision bounds for arbitrary varieties or motives.

\paragraph{Imported theorem and attribution.}
The geometric mechanism is classical. We use Asok--Dubouloz--\O stv\ae r,
\emph{Geometric models for algebraic suspensions}, arXiv:2112.08241v2,
Theorem~3.1.1, printed p.~8, and Remark~3.1.3, p.~9
(published in IMRN 2023, pp.~17788--17821).
For a smooth map to a contractible smooth base with a section, an
$\mathbb A^1$-equivalence off a smooth closed subscheme identifies the
total space with the cofiber of the induced map between the two normal
Thom spaces. We verify every hypothesis and this map below. Their
Construction~4.2.1 and Remark~4.3.2 already cover the general geometric
form; the new conclusion here is its precise application to the
parsimonious counting family and its complexity interpretation.
The proof of Theorem~3.1.1 was read in full. Its homotopy-purity input
is Theorem~2.3.1 there, extending Morel--Voevodsky,
\emph{$\mathbb A^1$-homotopy theory of schemes}, Section~3,
Theorem~3.23. This is motivic homotopy purity, distinct from an
inference from finite-field point counts.

\paragraph{Verification of the smooth map and normal framing.}
Put $B=\operatorname{Spec}\mathbb Z$, $W=\mathbb A^v_B$, and let
$\pi:X_\Phi\to W$ be projection to $x$. It admits the section
$s(x)=(x,0,0)$. For the relative equation
$F=t^2-t-\sum y_i f_i$, the partial derivatives in the fiber
coordinates are $2t-1,-f_1,\ldots,-f_h$. In the coordinate ring of $X$,
\[
                   (2t-1)^2-4\sum_i y_i f_i=1.
\]
Thus these derivatives generate the unit ideal. The usual smooth
hypersurface Jacobian criterion proves that $\pi$ is smooth of relative
dimension $h$ over the entire integral base, including characteristic
two. The same criterion is SGA~1, Expos\'e~II, Theorem~4.10(iii) and
Corollary~4.12, already checked for NN-C027.

Let $U=W\setminus Y_\Phi$. On $D(f_i)\subset U$ the equation solves
for $y_i$, leaving the independent fiber coordinates $t$ and the other
$h-1$ variables. Hence $\pi^{-1}(U)\to U$ is a Zariski locally trivial
$\mathbb A^h$-bundle and therefore an $\mathbb A^1$-weak equivalence.
The closed scheme $Y_\Phi$ is smooth over $B$, being a disjoint union of
integral sections. Its embedding in $W$ is a regular immersion of
codimension $v$.

The classes of $b_j=x_j^2-x_j$ frame its conormal bundle. Indeed, on
the component indexed by $a\in\{0,1\}^v$, their differentials are
\[
                  db_j=(2a_j-1)\,dx_j.
\]
Each coefficient is an integral unit. The conormal sequence for the
smooth closed pair therefore identifies these $v$ classes with a basis
of $\Omega^1_{W/B}|_Y$. Equivalently, on a neighborhood of each component
the ideal is generated by $x_j-a_j$ and $b_j=(x_j-a_j)(x_j+a_j-1)$,
whose second factors restrict to units. This also directly verifies
regularity and shows that the framing commutes with every base change.
No assertion that the whole ideal of an arbitrary finite union is
globally generated by these $v$ elements is used.

The fiber over $Y$ is exactly
\[
 X_Y=D_0\amalg D_1,\qquad D_e=\{x\in Y,t=e\}\simeq Y\times\mathbb A^h,
\]
because $t$ and $t-1$ are comaximal over $\mathbb Z$. Flat base change
along $\pi$ identifies its conormal bundle with the pullback of the
framed conormal bundle of $Y\subset W$. Thus both normal bundles have
compatible trivializations of rank $v$. The section $s|_Y$ lands in
$D_0$ as its zero section in the $y$ coordinates.

\paragraph{Computing the cofiber.}
The cited theorem gives the pointed cofiber sequence
\[
 \operatorname{Th}(N_{Y/W})\longrightarrow
 \operatorname{Th}(N_{X_Y/X})\longrightarrow (X,s(0)).
\]
Using the compatible frames, the first term is
$(\mathbb P^1)^{\wedge v}\wedge Y_+$ and the second is
$(\mathbb P^1)^{\wedge v}\wedge(D_0\amalg D_1)_+$.
Contraction of each $\mathbb A^h$ factor identifies this map with the
inclusion
\[
 \bigvee_{a\in\mathcal S}(\mathbb P^1)^{\wedge v}
 \longrightarrow
 \left(\bigvee_{a\in\mathcal S}(\mathbb P^1)^{\wedge v}\right)
 \vee\left(\bigvee_{a\in\mathcal S}(\mathbb P^1)^{\wedge v}\right)
\]
into its first summand. Its cofiber is the second summand. This proves
\eqref{eq:S017-motivic-wedge}, including $N=0$. The computation works
unchanged over every base ring, so no preservation of an unspecified
isomorphism under arbitrary derived base change is required.

\paragraph{Motives, higher cycles, and complexity.}
Applying the reduced motive functor sends a pointed wedge to a direct
sum and sends $(\mathbb P^1,\infty)^{\wedge v}$ to
$\mathbb Z(v)[2v]$. The chosen rational point splits off the unreduced
unit motive. This proves \eqref{eq:S017-motive}. Exact classical
interfaces are Mazza--Voevodsky--Weibel, \emph{Lecture notes on motivic
cohomology}, Example~13.17, Theorem~15.2 and Corollary~15.5 for the
Tate sphere and projective line, and Theorem~19.1 (printed p.~159) for
the higher-Chow comparison. Tate cancellation over perfect fields is
Remark~16.26, with the theorem credited there to Voevodsky's 2002
\emph{Cancellation theorem}. The representable cohomology formula is
simply the image of the finite wedge under the represented functor;
applying motivic cohomology and that comparison gives
\eqref{eq:S017-higher-chow}.

For $N>0$, the already established group
$\operatorname{CH}^v(X_{\Phi,k})=\mathbb Z^N$ proves that the reduced
motive is nonzero over every perfect field. Thus it also rules out
$\mathbb A^1$-contractibility. For an arbitrary field, base change to
its algebraic closure gives the same obstruction; over $\mathbb Z$,
base change to $\mathbb C$ suffices. For $N=0$ the wedge is empty, and
NN-S016 supplies in addition the actual integral affine-space
isomorphism. The source verifier has exactly $N=\#\operatorname{SAT}(\Phi)$
solutions with unique wire extension, so exact extraction of the
displayed multiplicity is parsimoniously \#P-complete, and its
vanishing is exactly UNSAT. These equivalences establish the asserted
coNP membership and hardness on this family. They make no claim that
general motivic isomorphisms have polynomial certificates.

\paragraph{Grothendieck dictionary and implications.}
The ingredients are smooth integral schemes, a framed regular closed
immersion, its normal cone, and the deformation-to-the-normal-cone
purity construction. The normal-cone/intersection-theory interface is
the same geometric framework as SGA~6, Expos\'e~VII, regular immersions;
the exact homotopy theorem used is the later theorem cited above, not
a theorem attributed to SGA~6. The resulting Tate decomposition is
stronger than the SGA~4/SGA~5 realization calculation in NN-C027.
It supplies actual integral motives after passage to a perfect field,
without a Tate conjecture or a deduction of motives from zeta functions.

Existence of a pure Tate realization, rationality, algebraicity of all
classes, and now an explicit unstable motivic sphere type do not supply
an efficient multiplicity extractor. E evaluates this multiplicity at
its source counting bound; it does not make it polynomial-time.
Neither A nor M is used. No canonical GCT multiplicity identification,
polynomial algorithm for SAT, or separation of VP and VNP follows.

